% ======================================================================
% Section 6 — Open problems
% ======================================================================

\section{Open problems}\label{sec:open}

The reformulation sharpens several questions into finite, explicit form.
We list them in what we judge to be increasing distance from the present
state of the framework.

\begin{enumerate}[leftmargin=1.7em,itemsep=4pt]

\item \textbf{The two hypotheses of Theorem~\ref{thm:SC}.} The paper's
  own frontier, and now the program's closest obligation. H1g is reduced
  to the \emph{strand census}: bound the number of arrangements and the
  phase freedom per arrangement, uniformly over the renormalization
  orbits (\S\ref{sec:scale-h1})---class-by-class work over resonance
  types, with the ratio-$2$ binary cascade the dominant first class and
  uniformity across classes itself a lemma. H2m's sharp form is the
  \emph{pair-line count}: the Lemma~\ref{lem:E} three-distance
  machinery applied to the junction draws, converting the cut law from
  hypothesis to theorem and pinning $\gamma_0$ (the chain-class bound of
  Theorem~\ref{thm:chaindilate} has the right shape with constants loose
  by orders of magnitude; the route is to count on the
  $\sum o=\mathrm{ov}$ slice). H0's route is the tangency-codimension
  argument. If all three close, Theorem~\ref{thm:SC} becomes
  unconditional and yields a uniform-in-$T$ closure of the $p^2$
  covering obstruction---hence, by the lossless equivalence, the
  conjecture at the covered rungs; the honest reading after five
  iterations of the obligation sequence (\S\ref{sec:scale-status}) is
  that the reformulation has a stubborn residual, and the theorem is
  stated so that the residual is exactly what these hypotheses name.

\item \textbf{The pair-selection law.} $(\mathrm{TAU}$-$_n)$ asserts that
  some pair certifies; the empirical law of \S\ref{sec:measure} says
  something stronger and unexplained: in every corpus measured, some pair
  certifies \emph{at a configuration that never covers}---a never-capable
  $(N,\text{order signature})$. A structural reason for the ubiquity of
  safe configurations would be the first general theorem of the program
  beyond the base lemmas, and the last ingredient needed to convert
  per-modulus verification into uniform proof on the unit side.

\item \textbf{The kernel world.} After the classifications, the entire
  residual of every corpus is one class: mixed order signatures at
  composite moduli (an order-$d$ residue contributes a bad set of size
  $N/d$, so moat counting has slack exactly there). The kernel world
  already contains the base rung's single open set,
  $V=(3,5,8,13)$, whose live pair obstruction is the covering of
  $\Z_{16}$ by the effective triple $(3,5,8)$. \emph{Post-submission
  (Remark~\ref{rem:postcomposite}):} the structure theory now exists in
  proved form --- the fibration theorem (kernel bad sets are pullbacks of
  unit bad sets at reduced moduli with the same threshold, so covering at
  composite $N$ transfers to the divisor lattice of $N$), together with
  the no-unit-assistance theorems in their provable part. The $p^2$ side
  is closed outright (the three-unit cell is empty at every prime; see
  the next item); the $pq$ side is now carried by the tensor translation
  of \S\ref{sec:pq} --- verified empty to $N\le143$ and reduced to the
  same input structure, a single unbounded input after this revision's
  discharge of Input E --- with the rigid-zone lifts the remaining open
  piece.

\item \textbf{The composite micro-cases, restated.} The original form of
  this problem --- ``no covering exists at $49$, $77$, $91$; the
  obstruction is unexplained'' --- is superseded
  (Remark~\ref{rem:postcomposite}): kernel-speed coverings exist there and
  are exactly the transferred rigid-zone configurations, and the unit
  safety of $77$ is proved. The corrected open problem is the \emph{no
  unit-assisted covering conjecture}: at $N=p^2$ and $N=pq$ with $p\ge5$,
  no covering family of at most four distinct bad sets exists in which the
  kernel subfamily does not already cover. Proved: families with at most
  two unit members (fiber-cap counting, with the sharp closed-form cap
  $2\lfloor s/6\rfloor+1+[s\equiv5\pmod6]$ for fibers of prime size $s$);
  three unit members (an inclusion--exclusion identity forced to an exact
  partition, killed by the Dirichlet floor $|A\cap\lambda A|\ge2$ on arc
  intersections); pure-unit families (reduced, via the $p$-multiples,
  to the rigid zone); and, post-submission, the entire $p^2$ side --- the
  $(1K,3U)$ cell is empty at every prime $p\ge5$ via the trichotomy
  lemmas (proved for $p\ge17$ resp.\ $p\ge31$ \emph{unconditionally}: the
  one-foot input is discharged by Lemma~\ref{lem:E}; finite-case-analyzed
  at the boundary primes; false and directly checked at $7,13$) plus the
  shear-rigidity theorem for $\pm$-colliding units
  (linear per-fiber drift against pinned covering configurations;
  verified at every prime $11\le p\le113$ with exhaustive ground truth at $p\le43$) --- and, in this revision, the $pq$ side of
  the two-unit and three-unit cells: the $(1,1,2)$ cell and both
  $(1K,3U)$ orientations are verified empty at every $N=pq\le143$ ($10$
  semiprimes, over a million family checks, margins $2$--$18$), and the
  row translation (Theorem~\ref{thm:R}) reduces their closure to the same
  remaining input, with the row census and the pinning of
  \S\ref{sec:pq-census}--\S\ref{sec:pq-verify} as the measured base.
  Open: the rigid-zone lifts, and the all-$p$ discharge of the single
  remaining named input. The boundary
  instance $N=9$ --- a genuine unit-assisted covering in exactly the
  closed $p^2$ cell shape --- shows the hypothesis $p\ge5$ is necessary.
  The named computational inputs are now precise
  objects (\S\ref{sec:gt-support}): the foot-count proposition \emph{in
  its consumed form} (the one-foot trichotomy at the two size pairs
  $(2k{+}2,2k{+}1)$ and $(2k{+}1,2k{+}1)$, viable differences confined to
  $\{2,(p-1)/2\}$) is \emph{proved} --- Lemma~\ref{lem:E}, the
  three-distance foot count, clean at every prime $p\ge31$ and
  cross-checked against the census to $p\le419$; the two-foot version is
  consumed by no current proof --- leaving the
  pinning-size lemma (relative-offset sets of at most $6$ elements with
  covering counts $28$ and $68$ independent of $p$, verified to
  $p\le113$, needed by the shear-rigidity closure, and now confirmed from
  the $pq$ side by Proposition~\ref{prop:pqpinned}) as the programme's
  single unbounded input, dischargeable not by
  bounded search but by the four-family covering classification
  (Conjecture~\ref{conj:SL}) --- which is the programme's route to an
  unconditional-in-$p$ statement.

\item \textbf{Rigid-zone termination.} Is the unit-capable set finite at
  each rung? Within the verified range (primes $\le 150$, composites
  $\le 111$) the $n=5$ zone is exactly $\{7,13,17,19,37\}$, but no
  finiteness argument is known, and the $n=4\to n=5$ transition
  ($11$ out, $\{17,19,37\}$ in) shows the zone is not monotone in the
  rung. Finiteness at one rung would make the classification program of
  \S\ref{sec:rigid} literally completable.

\item \textbf{The successor no-covering conjecture}
  (Conjecture~\ref{conj:successor}): for prime $N\ge 41$, four
  translates of $S_h$ never cover $\Z_{(N-1)/2}$. Counting never
  obstructs there ($4(2h+1)-3\ge N$), so the statement is purely
  multiplicative, the exact analogue one rung up of the $n=4$
  no-covering statement. The classifications of
  Theorems~\ref{thm:K4}--\ref{thm:K37} are the complete positive side of
  the ledger: every capable prime $\le 37$, classified.

\item \textbf{Is the decay fundamental?} The proved-coverage rate
  declines with range (\S\ref{sec:measure}): on the augmented battery
  the residual grows $6.56\%\to21.84\%$ over $V=24\to56$, with
  increments $+3.80$, $+5.36$, $+4.07$, $+2.06$ points per $8$ units of
  $V$---roughly linear growth with fluctuation, no established trend in
  either direction. The level fact is sharp: coverage fell below the
  program's $80\%$ continuation threshold at $V=56$ ($78.16\%$). The
  alternatives remain: either the decay is intrinsic to the
  reformulation (proved coverage has a range-dependent ceiling), or it
  is an artifact of the battery---specifically of the absence of a
  kernel-world lemma, since the kernel world is $100\%$ of the residual.
  Five data points do not decide the question; a kernel-world structure
  theorem would.

\item \textbf{Plain-regime growth.} On the any-pair basis the plain share
  is stable across rung and range ($38.3\%$ at $n=4$; $34.3\%$--$36.0\%$
  at $n=5$), but the verified-only sets are all plain, and the
  polygon lemma's share shrinks structurally as $T$ grows (a residue is
  polygon-certifiable only when some $j\mid N$ with $j\le T$ avoids all
  residues). Does the plain share grow with $n$, and at what rate?

\item \textbf{Higher rungs.} At $n=6$ ($T=7$, five bad sets per pair) the
  rigid zone should be denser (more translates available) and the kernel
  world larger. The program's method---run the battery, map the new
  rigid zone, generate the classifications the new level demands, replace
  the ports that die---is exactly what was executed at $n=5$; whether it
  survives a third rung is untested. \emph{Post-submission}
  (\S\ref{sec:higherrung}): the census+shear method \emph{does} survive
  the third and fourth rungs. At the critical cell of $T=8$ the fiber
  census is exactly eleven families with placement counts constant in
  $p$ (Proposition~\ref{prop:T8census}, seventeen primes); the
  $\pm$-colliding-only law holds at both rungs in the mechanism range;
  the drift law is $T$-free (Lemma~\ref{lem:rungT}); the shear step
  survives the growing pinned sets because the growth is dimensional
  (solution sets $\approx k^2$, projections saturating at $42$ against
  $|U|=(T-2)k+\epsilon-1$, exit time $\le2$); and the frontier cells are
  empty --- $(1K,4U)$ at $T=7$ for every prime $11\le p\le61$, $(2K,4U)$
  at $T=8$ for every prime $17\le p\le43$ tested. The rigid-zone lift
  itself (pure-unit families) remains open at every rung, as does the
  $T=8$ $(1K,5U)$ cell, whose five-difference census is a zoo (840
  $\pm$-distinct families at $k=2$) --- the honest boundary of the port,
  and now the object of the conditional scaling-closure attack of
  \S\ref{sec:scale}: Theorem~\ref{thm:SC} closes it under H1g and the
  H2 pair, with the strand census as the named remaining obligation.

\item \textbf{Bridges to neighboring methods.} Three external toolkits
  interface naturally with the reformulation's finite statements.
  (a)~The covering lemmas now have a \emph{proved} bridge: the group-ring
  translation (\S\ref{sec:groupring}) turns them into sign conditions on
  a sparse identity in $\Z[\Z_p]$, proved as Theorem~\ref{thm:GT} and
  machine-verified on every covering in the $p\le61$ census; its census
  law (four difference families with $p$-independent placement counts,
  Conjecture~\ref{conj:SL}) is the object the remaining two lemmas must
  prove, and it sits precisely in the Lam--Leung--de~Bruijn--Schoenberg
  theory of sparse cyclotomic divisibility.
  (b)~The no-covering
  statements (Conjecture~\ref{conj:successor} and its $n=4$ ancestor)
  concern coverings of cyclic groups by \emph{multiplicative
  arcs}---images of intervals under multiplication. Additive-combinatorial
  covering theory (Kneser-type theorems, Cauchy--Davenport and weighted
  variants) is the classical toolkit for covering questions in $\Z_N$;
  whether any of its forms applies to multiplicative arcs is open, and a
  positive answer at prime $N$ would settle the successor conjecture.
  (c)~Covering-radius problems for lattice zonotopes are governed
  matroidally, and log-concave polynomial methods have turned matroidal
  counting statements into general theorems (Mason's ultra-log-concavity
  conjecture being the template); whether the pair-sum covering problems
  admit a zonotope--matroid encoding whose associated polynomial has a
  stability property is the corresponding question here---if they do, the
  kernel world would be its near-equality cases, and the sign-mixed
  sparse divisibility of \S\ref{sec:gt-outlook} is the natural first
  candidate. The first bridge is exploited; the other two are recorded as
  directions, not claims. \emph{Post-submission} (final revision): the
  third bridge was subsequently engaged as a full strand and is now
  closed. The zonotope--Ehrhart encoding exists---the Lonely-Runner
  zonotope of the instance at denominator $N$, its dilate family
  $mZ(V)$---and its arithmetic is entirely stable: the $h^*$-polynomials
  of both canonical families (the harmonic instances $v=(1,\dots,n)$ and
  the rung-$2$ instances $v=(1,\dots,n{-}1,2n)$) are real-rooted,
  log-concave, and Newton at every $n\le 10$, verified in exact
  rational arithmetic at $45$ digits \cite{companion-paper}. The
  geometric control the bridge asked for, however, is false in its
  natural form: the deepest-hole reading---the covering optimum
  attained at the zonotope's center, $\rho=\delta$---fails from $n=5$
  onward with exact rational witnesses (at the harmonic $n=5$ instance
  $m=5$, $\rho=16/47>\tfrac13$ \emph{exactly}, closed in the companion
  paper's fourth revision by a balance-family witness and two new box
  certificates---the search record
  $97/288\to607/1792\to4183/12288$ retained as provenance; the two
  $n=6$ boundary rows resolved as strict refutations, $17/47$ and
  $455/1269$), the modular law that would
  govern it (deepest hole at the center exactly when $n\mid m$) breaks
  at $n=5$ (the odd multiples $5$ and $15$ fail exactly; the even and
  higher multiples $10,20,25,30$ are certified, $m=10$ closing with
  the companion paper's third revision) and is extinct at $n=6$, and
  the strictness of $h^*$ neither
  tracks nor predicts the covering gap (rank correlation at the noise
  level). The bridge's question is thus answered in the negative for
  the natural encoding, and with the same shape as the internal
  program's boundary: arithmetic fidelity, no margin. The witnesses,
  the exact certification ladder, and the $h^*$ addendum are the
  subject of the companion paper \cite{companion-paper}.
\end{enumerate}
